\documentclass[11pt]{article}
\usepackage[margin=1in]{geometry}
\usepackage[T1]{fontenc}
\usepackage{lmodern,microtype,amsmath,amssymb,amsthm,booktabs,needspace}
\usepackage[colorlinks=true,linkcolor=blue,urlcolor=blue]{hyperref}
\setlength{\parindent}{0pt}
\setlength{\parskip}{3pt}
\newtheorem{proposition}{Proposition}
\title{Can cash transfers improve the housing allocation?}
\author{Discussion note for Tommaso De Santo}
\date{September 5, 2026}
\begin{document}
\begin{center}
{\LARGE Can cash transfers improve the housing allocation?}\par
\vspace{0.6em}
Discussion note for Tommaso De Santo \quad September 5, 2026
\end{center}
\vspace{0.4em}

\textbf{Answer.} Yes, if the planner can commit to household-specific cash
transfers when households are young and when they are old. In the stationary
case described below, a small reform moves housing from old owners to young
owners at every date. All household cohorts can retain their lifetime utility,
and a resource surplus remains after accounting for every initial housing
title. Prices and private choices clear every housing market along the path.

The permission matters. A promise to tax a household when old can support a
transfer before its housing purchase. This provides public financing across
ages, even though the private mortgage rule remains unchanged. The result
establishes constrained inefficiency relative to this committed transfer class.
It does not establish the same claim for a single round of cash gifts, or
identify a welfare loss caused solely by households ignoring their effects on
prices. No fertility externality is used.

\section{The comparison}

Use the existing household preferences and notation. A young household has
adult consumption $x=c-\chi n$ and adult space $s=h-\kappa n$, with utility
$\log x+\alpha\log s+\vartheta\log n$. Old utility is
$\log c^2+\gamma\log h^2+\omega_B\log e$. Hold each household's fertility
at its original choice and keep cohort sizes fixed. The exercise therefore
compares housing use and the associated consumption and estates.

Start at a stationary equilibrium with mass $m>0$ of owners and mass $m_R>0$
of renters at each age. Within each tenure group, households have common
income, liquid wealth, and fertility; their original ownership tastes can
differ. Young owners have a strictly binding down-payment limit, positive
saving, and a slack physical size cap. Old owners have slack retention and
estate-composition bounds. Young and old renters are at their rental size cap.
Tenure remains optimal in a sufficiently small neighborhood, as discussed below.
This is a particular stationary case of the model, not a theorem for every
entrant distribution or every set of binding constraints.

Write $u=q r=(1+q\tau^p-q)P$ for the stationary cost of housing services.
Use $h^Y,h^O,c^O$ for stationary young-owner housing, old-owner housing, and
old-owner consumption. Define three preference sums and the young owner's
housing value in units of consumption:
\begin{equation}
 A=1+\gamma+\omega_B,\quad
 \rho=1+\beta A,\quad B=1+\beta(1+\omega_B),\qquad
 MV^Y=\frac{\alpha x}{s},\quad s=h^Y-\kappa n.
\end{equation}
The original saving and old-housing conditions give
\begin{equation}
 c^O=\frac{\beta x}{q},\qquad
 h^O=\frac{\gamma\beta x}{q u},\qquad MV^O=u.
 \label{eq:ss}
\end{equation}
Thus $MV^Y>u$ is the housing-value gap generated by the binding purchase limit.

\textbf{Transfers and ownership.} A young household receives cash $\ell_t^Y$
before closing, and a previously announced payment $\ell_{t+1}^O$ when old.
Negative payments are enforceable lump-sum taxes. A household receives its
announced payments under either tenure choice; they are not housing subsidies.
The young-owner limit remains
$(1-\phi)P_t h_t^Y\leq b+\ell_t^Y$. Private saving remains nonnegative.
The ordinary property-tax rebate $T_t$ retains its original timing and budget.

We must also identify the owners of the initial titles not held by current old
households. For this proposition, assign the remaining
$R=\bar H-mh^Y$ titles to a passive outside owner. This owner has no domestic
housing demand, sells its initial titles at the reform price, and can trade the
world bond subject to its present-value budget. It values additional goods and
receives the residual of the cash-transfer account at each date. Its title
revaluation and its entire transfer stream enter the welfare comparison.
This is an explicit completion of the original model's unspecified initial
ownership, not a claim that zero intermediary profits eliminate initial owners.
The planner can enforce the announced payments to and from this owner.
With sufficiently small reforms and slack outside credit capacity, all promised
payments are feasible without a Ponzi scheme.

\section{A fully clearing improving path}

\Needspace{12\baselineskip}
\begin{proposition}[Committed cash transfers]
Under the preceding assumptions, suppose $MV^Y>u$ and
\begin{equation}
 \lambda=\frac{\alpha h^O}{\rho s}<1.
 \label{eq:contraction}
\end{equation}
Then a sufficiently small committed cash-transfer reform is a Pareto
improvement. Young owners occupy more housing and old owners less at every
finite date. All household cohorts can be indifferent and the passive initial
title owner strictly better off. Every housing market and each date's transfer
budget balance. The allocation converges to the original steady state.
\end{proposition}

The convergence condition limits how strongly a small consumption change in
one cohort changes the next cohort's compensated housing choice. It is a
sufficient condition for this construction, not a claim about every policy
transition in the model.

\textbf{First date.} Let $\epsilon>0$ be a small proportional increase in the
current service cost. Set
\begin{equation}
 P_0=P+\frac{u\epsilon}{1+q\tau^p},\qquad
 P_t=P\quad(t\geq1).
\end{equation}
This implies $u_0=u(1+\epsilon)$ and $u_t=u$ thereafter. The price path is
supported by the transfers below; it is not a price control.
Let $\eta=(1+\omega_B)/A$. Exact compensation of the initial old owners gives
\begin{equation}
 c_0^O=c^O(1+\epsilon)^{\gamma/A},\qquad
 h_0^O=h^O(1+\epsilon)^{-\eta},\qquad e_0^O=\omega_B c_0^O/q.
 \label{eq:initialold}
\end{equation}
These are their optimal private choices at the new user cost. Their utility
is unchanged, and they release housing. Keep renter consumption and housing
unchanged. Market clearing then requires
\begin{equation}
 h_0^Y=h^Y+h^O-h_0^O,\qquad
 x_0=x\left(\frac{s}{h_0^Y-\kappa n}\right)^{\alpha/\rho}.
 \label{eq:firstyoung}
\end{equation}
The young owner has more space and less adult consumption, with the same
lifetime utility. Its old-age consumption and estate adjust through saving.

\textbf{Every later date.} Define $X_t=x_t/x$. The preceding young cohort
chooses old housing $h^OX_{t-1}$ at the unchanged future user cost. Clearing
housing and preserving the next young cohort's lifetime utility give
\begin{equation}
 h_t^Y=h^Y+h^O(1-X_{t-1}),\qquad
 X_t=\left[1+\frac{h^O}{s}(1-X_{t-1})\right]^{-\alpha/\rho},
 \quad t\geq1.
 \label{eq:recursion}
\end{equation}
Since $X_0<1$ and the derivative of this map is at most $\lambda<1$ on
$X\leq1$, the sequence increases to one. Young housing is above its original
level at every date, old housing is below it, and both converge to baseline.
All sufficiently small changes preserve the strict private constraints stated
above. This recursion verifies the infinite continuation rather than fixing
future prices without clearing future markets.

\textbf{Private implementation.} Pay each young owner the closing cash needed
for its chosen home. Choose its committed old-age payment from its original
lifetime budget:
\begin{align}
 \ell_t^Y&=(1-\phi)P_t h_t^Y-b,\label{eq:younggrant}\\
 \ell_t^Y+q\ell_{t+1}^O
 &=\rho x_t+u_t h_t^Y+\chi n-y-b-T_t-qT_{t+1}.
 \label{eq:lifetimegrant}
\end{align}
The saving Euler equation, old demand conditions, and binding purchase limit
then support the proposed allocation. A small change leaves the young
housing multiplier positive. Cash payments to initial old households follow
from their original dated budgets. Initial old renters and young renters
receive $(u_0-u)h_R^{\max}-(T_0-T)$; later renters need no transfer. Each date,
the passive owner receives the negative of the other groups' aggregate net
transfers, so the government issues no debt.

The passive owner can have payments of either sign across dates. Its access to
the world bond, together with enforceable future household taxes, is the
financial capacity that supports this construction. Date-by-date government
balance alone would not prove that this financing is feasible.

\section{The resource gain and its interpretation}

Let $\Delta W$ be the change in the passive owner's present-value resources,
including $R(P_0-P)$ and its complete transfer stream. Summing the original
dated budgets, using housing clearing and the property-tax rebate rule, gives
the exact identity
\begin{equation}
 \Delta W=-m\left[(1+\omega_B)(c_0^O-c^O)
       +B\sum_{t=0}^{\infty}q^t(x_t-x)\right].
 \label{eq:resources}
\end{equation}
Every initial title is counted. House-price revaluations and domestic taxes
cancel after consolidation; changes in consumption and warm-glow estates
remain. The negative consumption changes of successive young cohorts can
pay for initial old-owner compensation.

Differentiate the constructed path at $\epsilon=0$. Equations
\eqref{eq:initialold}--\eqref{eq:recursion} imply
$\mathrm dh_t^Y/\mathrm d\epsilon=h^O\eta\lambda^t$ and
$\mathrm dx_t/\mathrm d\epsilon=-(MV^Y/\rho)h^O\eta\lambda^t$.
Substituting these in \eqref{eq:resources} and using \eqref{eq:ss} yields
\begin{equation}
 \left.\frac{\mathrm d\Delta W}{\mathrm d\epsilon}\right|_{0}
 =\frac{m h^O\eta}{1-q\lambda}\,(MV^Y-u)>0.
 \label{eq:gain}
\end{equation}
The gain has the desired source: the housing is worth more to the young than
to the old. All household utilities have already been compensated, so the
remaining positive resources establish a strict Pareto improvement. The proof
assigns the surplus to the passive owner for convenience; no welfare weights
or fertility benefits enter the argument.

The initial young owners can also gain strictly. At a fixed small reform,
raise $x_0$ slightly above \eqref{eq:firstyoung}, keeping $x_0<x$, and use the
same recursion from date one onward. Their housing is unchanged and their
lifetime utility rises. By continuity, some positive resource surplus remains,
all later cohorts retain their utility, and housing still moves toward the young.

\Needspace{7\baselineskip}
\textbf{Real transfer costs.} The proposition uses zero additional real cost.
An explicit extension charges $C(h_t^Y-h^Y)$ goods per affected owner at date
$t$, with $C(0)=0$ and $C'(0)=K$. Deduct these costs from the passive owner's
cash account. The derivative in \eqref{eq:gain} then replaces $MV^Y-u$ with
$MV^Y-u-K$. This is a cost of the extra housing reassigned at each date, not
automatically a fixed moving cost or a cost of the change from the previous date.

\textbf{Tenure choice.} The transfers are fixed by household identity before
the new choice. Holding fertility fixed means an original owner considering
renting must use its own $n$, and conversely. If the two tenures originally
choose different fertility levels, a household at the original taste cutoff
can have a strict preference for its original tenure in this comparison.
A uniform positive margin and continuity preserve tenure for a sufficiently
small reform, including a continuous taste distribution. This condition must
be checked rather than inferred from a fixed owner share.

\section{Why a one-time gift is a different question}

A transfer only at the initial date cannot generally supply the compensation
in \eqref{eq:lifetimegrant} for all later cohorts. Their housing demand changes
as the initial cohort ages, so prices adjust again. In the illustrative
stationary case, the linearized equilibrium after transfers cease has one
converging price root, $-0.05518$, and one divergent root, $3.30801$. The
converging price response alternates in sign. With renters at their housing
caps, their lifetime welfare also alternates in sign, so a nonzero small
converging price response hurts some future cohorts.

The separate review derives the nonlinear local argument and treats the
zero-price-response case using the complete resource account. Its conclusion
is a local obstruction in this stationary regime. It is not a global
impossibility theorem for one-time transfers under all preferences, initial
conditions, or tenure responses. The derivative roots alone would not establish
that stronger claim. The earlier fixed-future-price, four-household examples
omit both this continuation and some initial title owners; they cannot settle
the complete-economy comparison.

\textbf{Recommendation.} Use the committed-transfer result if we accept public
financing across ages as part of the constrained planner's powers. State that
permission before the proposition. If the intended planner is restricted to
one initial cash redistribution, the present model does not yet supply the
desired general result, and the local obstruction is substantive. No new
fertility motive is needed to obtain the positive theorem above.

\section{Verification}

Independent derivations, original-budget accounting, and twelve direct
household optimizations support the proof. Accounting errors are below
$9\times10^{-16}$; optimized choices match within $1.8\times10^{-7}$.
The minimum tenure margin exceeds $0.0357$ and $\lambda=0.26088$.
These are illustrative checks, not calibration evidence. The reviews and
receipts are indexed in \path{output/model/simplified_olg_amendments/README.md}.
Reproduce the checks, including independent optimizations with SciPy, with
\par\smallskip
{\footnotesize
\texttt{python3}\ \path{code/model/tools/verify_simplified_olg_planner_benchmarks.py}\ \path{--original-optimizers}\par}
The author-controlled manuscript and quantitative model were not edited.
\end{document}
